% answers_short.tex —— 2001 年·春季上海卷 · 简短答案速查
% 由 tools/make_short_answers.py 自 answers.tex 的 \gksolhead/\gkanswer 提取，
% 并与 2001/2001_数学真题_春季上海_答案要点.md 逐题比对后生成，勿手改（改源 MD 后重跑）。

\begin{gknotice}
\textbf{使用说明}：本节只给最终答案便于快速核对，完整推导见解析版。答案为编写者独立推导，\textbf{非官方标准答案}。
\end{gknotice}

\gkpart{填空题}
\begin{itemize}[leftmargin=2.6em]
  \item[\textbf{1.}] $f^{-1}(x)=-\sqrt{x-1}$（$x\ge1$）。
  \item[\textbf{2.}] $z=1+i$。
  \item[\textbf{3.}] $2\pi$。
  \item[\textbf{4.}] $20$。
  \item[\textbf{5.}] $\dfrac{x^2}{9}-\dfrac{y^2}{16}=1$。
  \item[\textbf{6.}] $(x-1)^2+(y-1)^2=1$。
  \item[\textbf{7.}] $e^2$。
  \item[\textbf{8.}] $90^\circ$。
  \item[\textbf{9.}] $\dfrac45$。
  \item[\textbf{10.}] $a\ast b+c=b\ast a+c$。
  \item[\textbf{11.}] ①，$0$（也可填④，$\dfrac\pi2$）。
  \item[\textbf{12.}] $219.01$ 元，甲比乙多。
\end{itemize}

\gkpart{选择题}
\begin{center}\begin{tabular}{ccccc}
\toprule
题号 & 13 & 14 & 15 & 16 \\
\midrule
答案 & C & C & D & B \\
\bottomrule
\end{tabular}\end{center}

\gkpart{解答题}
\begin{itemize}[leftmargin=2.6em]
  \item[\textbf{17.}] $\left(-2,-1\right)\cup\{3\}$。
  \item[\textbf{18.}] $\sqrt{1-k}$。
  \item[\textbf{19.}] (1) $a=\dfrac1{\sqrt{1+h^2}}$；(2) $h=1$ 时 $V$ 最大，最大值为 $\dfrac16$ 立方米
  \item[\textbf{20.}] (1) 见解析；(2) $\arccos\dfrac{12\sqrt2}{25}$
  \item[\textbf{21.}] (1) $x^2+\dfrac{y^2}{2}-ax-\dfrac{by}{2}=0$；(2) 交点集合为 $(0,0)$、$(a,0)$、$(0,b)$；按三点是否重合分情况计数
  \item[\textbf{22.}] (1) $S_{n+1}=\dfrac{S_n+4}{2}$；(2) 不存在这样的自然数 $c$ 和 $k$
\end{itemize}
